\chapter{フリードマン膨張、放射、暗黒物質}

\section{基本的宇宙論方程式の復習}

銀河団より大きなスケールでは、宇宙は一様であるとみなされる。座標を選んで我々が\(x=0\)に位置するようにし、代表的な銀河、あるいは単にハッブル流れとともに運動する想像上の粒子を\(x=1\)に置く。すると、我々からのその物理的距離は、定義により、スケール因子\(a(t)\)である。\(x=2\)にある点は距離\(2a(t)\)にあり、以下同様である。

\begin{figure}[ht]
\centering
\includegraphics[width=0.78\textwidth]{lecture_03_figure_02.png}
\caption{\(x=0\)を中心とする一様な球、\(x=1\)にある共動銀河、および\(a(t)\)に対するニュートン的なエネルギー釣り合い。}
\end{figure}

半径\(a(t)\)の球の内部の物質だけが、\(x=1\)にある粒子への重力に寄与する。その粒子に単位質量を与える。その運動エネルギーは
\begin{equation}
\frac{1}{2}\dot a^2.
\end{equation}
球に囲まれた質量は
\begin{equation}
M_{\mathrm{enc}}=\frac{4\pi}{3}a^3\rho,
\end{equation}
であるから、重力ポテンシャルエネルギーは
\begin{equation}
-\frac{G M_{\mathrm{enc}}}{a}
=
-\frac{4\pi G}{3}a^2\rho.
\end{equation}
\editorialnote{講義のタイムスタンプ00:05:30において、板書には運動項と密度に依存するポテンシャル項が明瞭に示されている。ここに書かれた完全な係数は、口頭での導出から従う。}

エネルギー保存により
\begin{equation}
\frac{1}{2}\dot a^2-\frac{4\pi G}{3}a^2\rho=\text{const}.
\end{equation}
この定数は\(-k\)と書かれ、標準的な規格化が実際にその記法を確定するのは方程式を\(2\)で掛けた後だけであるという、いつもの注意が付く。標準的な形では、
\begin{equation}
\dot a^2=\frac{8\pi G}{3}a^2\rho-k,
\end{equation}
であり、\(a^2\)で割ると
\begin{equation}
\left(\frac{\dot a}{a}\right)^2=\frac{8\pi G}{3}\rho-\frac{k}{a^2}.
\end{equation}
が得られる。これが、議論全体で用いられる形のフリードマン方程式である。

\begin{figure}[ht]
\centering
\includegraphics[width=0.56\textwidth]{lecture_03_figure_03.png}
\caption{フリードマン方程式において孤立させた曲率の寄与\(-k/a^2\)。}
\end{figure}

曲率項は経験的に非常にゼロに近く、この水準では単に無視することが非常に良い近似となるほどに近い。それは確かに空間の曲率に対応するが、観測的にはここではその符号すら決定できないほど小さい。

同じ方程式は一般相対論からも得られる。概略的な形では、時間-時間アインシュタイン方程式は
\begin{equation}
R_{00}-\frac{1}{2}g_{00}R \propto T_{00}.
\end{equation}
\editorialnote{ここでの議論はテンソル構造のみを与え、正確な全体の数値係数は与えない。}
膨張する一様な計量では、幾何学側は\(\dot a/a\)と\(k/a^2\)の組み合わせを生み、物質側はエネルギー密度を生む。この段階では、ニュートン的な議論とアインシュタイン方程式は同じことを言っている。

それは\(\rho\)の意味も変える。ここから先は、\(\rho\)をエネルギー密度と呼ぶ方がよい。共動系でゆっくり運動する物質に対しては、それは単に静止エネルギー密度である。

\section{物質支配}

ここで、エネルギー密度が共動系で静止している物質によって担われていると仮定する。一辺が座標1単位の立方体を取る。その物理的な辺の長さは\(a(t)\)であるから、その物理的な体積は\(a^3(t)\)である。その共動立方体の内部の質量エネルギーが\(M\)であり、物質が境界を横切るのではなく流れとともに運動するならば、
\begin{equation}
\rho_m=\frac{M}{a^3}.
\end{equation}

これが物質支配の場合である。この理想化では粒子は箱の中で静止しているので、温度はゼロとみなされ、圧力もゼロとみなされる。圧力には、粒子が箱の境界を横切ることが必要であろう。

\(k\approx 0\)とすると、方程式は
\begin{equation}
\left(\frac{\dot a}{a}\right)^2=\frac{8\pi G}{3}\frac{M}{a^3}.
\end{equation}
となる。冪乗則を試すと、
\begin{equation}
a(t)=Ct^p.
\end{equation}
すると
\begin{equation}
\frac{\dot a}{a}=\frac{p}{t},
\end{equation}
であるから、左辺は\(t^{-2}\)のようにスケールし、右辺は\(t^{-3p}\)のようにスケールする。冪を合わせると
\begin{equation}
3p=2,
\qquad
p=\frac{2}{3}.
\end{equation}
が得られる。したがって
\begin{equation}
a(t)\propto t^{2/3}.
\end{equation}
係数を保つと
\begin{equation}
C^3=6\pi G M.
\end{equation}

この解の要点は冪乗則である。物質支配の宇宙は時間の3分の2乗のように膨張し、線形ではない。この法則は、ビッグバンからおよそ\(10^5\)年後から数十億年後までの長い中間時代を記述する。現在の宇宙は文字通りには物質支配ではない。なぜなら今は暗黒エネルギーが支配しているからであるが、通常の物質と宇宙論的歴史の大きな部分はこのように記述される。

\section{放射支配}

それより前は、エネルギー密度は放射が支配していた。理由は単純である。宇宙には膨大な数の光子が含まれており、陽子一つあたりおよそ\(10^9\)個の光子が存在し、宇宙が圧縮されるとそれらの光子はより高いエネルギーを持つようになる。

光子には静止質量はないが、エネルギーはある。そのエネルギーは
\begin{equation}
E_\gamma=h\nu=\frac{hc}{\lambda}.
\end{equation}
である。決定的な仮定は、宇宙が膨張するにつれて光子の波長もそれとともに伸びるということである：
\begin{equation}
\lambda\propto a.
\end{equation}
箱の中の波が正しい描像を与える。箱が大きくなれば、モードの波長は箱とともに大きくなる。同じ考えがここでは膨張する宇宙に対して用いられている。

\(\lambda\propto a\) であるから、各光子のエネルギーは次のようにスケールする。
\begin{equation}
E_\gamma\propto a^{-1}.
\end{equation}
再び共動単位体積を考える。光子は出入りするが、平均すると内部の数は一定に保たれる。すると、箱の中の全光子エネルギーも \(a^{-1}\) のようにスケールする。これを物理的体積 \(a^3\) で割ると、放射エネルギー密度は次のようにスケールする。
\begin{equation}
\rho_r\propto a^{-4}.
\end{equation}

その余分な因子 \(a^{-1}\) がすべての違いである。物質は宇宙が膨張するために薄まる。放射はその理由に加えて、各光子が赤方偏移するために薄まる。

したがって、放射優勢宇宙に対するフリードマン方程式は次のようになる。
\begin{equation}
\left(\frac{\dot a}{a}\right)^2=\frac{8\pi G}{3}\frac{M_r}{a^4}.
\end{equation}
\editorialnote{口頭の議論では、この新しい定数に記号 \(M\) を再利用している。ここで添字 \(r\) を加えたのは、物質の場合との混同を避けるためだけである。}

同じ試行形 \(a(t)=Ct^p\) を用いる。左辺は再び \(t^{-2}\) のようにスケールし、右辺は \(t^{-4p}\) のようにスケールする。べきを一致させると
\begin{equation}
4p=2,
\qquad
p=\frac{1}{2}.
\end{equation}
したがって
\begin{equation}
a(t)\propto t^{1/2}.
\end{equation}
係数を保つと
\begin{equation}
C^4=\frac{32\pi G}{3}M_r.
\end{equation}

\begin{classroomqa}
\lecturetimestamp{00:35:41}
\audiencequestion{なぜそこに \(G\) があるのか。}
\lecturerresponse{重力を及ぼすのは静止質量だけではなくエネルギー密度だからである。ゆっくり動く物質ではこの区別を見落としやすいが、光子はそれを我々に強いる。光子は静止質量を持たないが、それでも重力を及ぼす。}
\end{classroomqa}

放射で満たされた箱はこの点を具体的にする。光子で満たされた箱を押す。箱が加速すると、光子に向かって動く壁は光子から遠ざかる壁よりも強く打たれるので、放射圧に対してより強く押さなければならない。したがって箱は余分な慣性を持ち、等価原理によりその余分な慣性は余分な重力質量も意味する。これが、エネルギーが重力を及ぼすと言うためにここで用いられる光子の箱の議論である。

\begin{classroomqa}
\lecturetimestamp{00:38:11}
\audiencequestion{これらの方程式のいずれかは、時間とともに減少する \(a\) で解けるか。}
\lecturerresponse{はい。曲率項 \(k/a^2\) を復元すれば、膨張した後に収縮する解がある。それは我々の宇宙がしているように見えることではないが、方程式はそれを許す。}
\end{classroomqa}

\begin{classroomqa}
\lecturetimestamp{00:38:44}
\audiencequestion{エネルギー保存はどうなったのか。}
\lecturerresponse{\(k\) 項を復元し、方程式に \(a^2\) を掛け戻すと、宇宙論的方程式はまさに保存則である。片側の物質または放射は簿記の一部にすぎない。膨張が運動学的部分と重力的部分を寄与し、全体の収支は保存される。}
\end{classroomqa}

\section{物質と放射の共存}

実際の宇宙論には物質と放射の両方が含まれる。それらの密度は異なる仕方でスケールする。
\begin{equation}
\rho_m\propto a^{-3},
\qquad
\rho_r\propto a^{-4}.
\end{equation}
したがって、小さい \(a\) では放射が優勢である。大きい \(a\) では物質が優勢である。時間を遡ると、初期宇宙はしたがって放射優勢であったに違いない。

\begin{figure}[ht]
\centering
\includegraphics[width=0.68\textwidth]{lecture_03_figure_05.png}
\caption{スケール因子の関数として物質と放射のエネルギー密度を比較するための軸。}
\end{figure}

\editorialnote{講義タイムスタンプ 00:44:45 において、黒板には比較プロットの始めの軸だけが示されている。その直後に議論される物質と放射の曲線は、ここに描き直すのではなく本文で記述されている。}

物質の曲線は \(a^{-3}\) のように減少する。放射の曲線は \(a^{-4}\) のように減少する。両方とも \(a\to 0\) で発散するが、放射の曲線はそこでより速く上昇するので、二つの曲線は等しいスケール因子で交差する。粗い見積もりでは、この転移はおよそ \(10^4\) 年の後、おそらくもう少し後になる。

同じ点は、\(a(t)\) 対 \(t\) のプロットで表すこともできる。初期には、膨張は \(t^{1/2}\) のように振る舞う。後には \(t^{2/3}\) のように振る舞う。変化は跳躍ではなく、比較的狭い時間範囲を通る連続的なクロスオーバーである。

\begin{classroomqa}
\lecturetimestamp{00:44:04}
\audiencequestion{宇宙が年を取るにつれて、べき \(p\) は \(1/2\) から \(2/3\) へ連続的に移るのか。}
\lecturerresponse{はい。跳躍はしない。変化は連続的であるが、かなり狭い時間範囲で起こる。}
\end{classroomqa}

\section{温度、再結合、最終散乱}

次の量は温度である。典型的な光子のエネルギーが \(a^{-1}\) のようにスケールするなら、放射浴の温度も同様にスケールする。
\begin{equation}
T\propto a^{-1}.
\end{equation}
\(a\) が2倍になれば、温度は半分に切られる。

現在の光子温度をおよそ \(3\,\mathrm K\)、電離温度をおよそ \(3000\,\mathrm K\)、すなわちほぼ太陽表面の温度とする。すると
\begin{equation}
\frac{a_{\mathrm{today}}}{a_{\mathrm{ion}}}\sim 10^3.
\end{equation}
後期の物質優勢則 \(a\propto t^{2/3}\) を用いると、
\begin{equation}
\left(\frac{t_{\mathrm{today}}}{t_{\mathrm{ion}}}\right)^{2/3}\sim 10^3,
\end{equation}
したがって
\begin{equation}
\frac{t_{\mathrm{today}}}{t_{\mathrm{ion}}}\sim 1000^{3/2}\sim 3\times 10^4.
\end{equation}
この見積もりの水準で宇宙の年齢を \(10^{10}\) 年とすると、
\begin{equation}
t_{\mathrm{ion}}\sim 3\times 10^5\ \mathrm{years}.
\end{equation}

これは水素とヘリウムが完全に電離した状態でなくなる時代である。それより前では、宇宙は高温のプラズマである。なぜプラズマは不透明なのか。自由電荷が電磁波に強く応答するからである。電場が電荷を前後に駆動し、媒質は導体のように振る舞い、導体は不透明である。これが太陽表面との関連する類推である。

宇宙が十分に冷えて電子が原子核と再結合できるようになると、通常の中性原子が形成され、宇宙ははるかに透明になる。この遷移は無限に鋭いわけではない。それはぼやけた場所、あるいはぼやけた時間間隔であるが、宇宙論にとって十分に鋭いので、最後の散乱面と呼ばれている。

外を見ることは時間を遡って見ることを意味する。近くの銀河は最近の過去を示す。より遠い銀河はより早い時代を示す。十分遠くでは、どの視線も宇宙がまだ電離して不透明であった時代に到達する。電磁放射はそれより前のことを示すことができない。

他のメッセンジャーは原理的にはより早い時代を探ることができる。ニュートリノは電気的に中性であり、より早く脱結合するであろう。重力子は、もし観測できたなら、さらに早い時代から来るであろう。しかし、それらが初期のプラズマを通り抜けることを可能にする同じ相互作用の弱さが、それらを極めて検出しにくくもする。

もし宇宙がかつてあらゆる方向で太陽表面と同じくらい高温であったなら、なぜ空は目がくらむほど明るくないのか。最後の散乱面が後退しており、そこから届く光子が\(3000\,\mathrm K\)を\(3\,\mathrm K\)から分けるのとほぼ同じ\(10^3\)の因子で赤方偏移しているからである。我々に届くのはマイクロ波放射であり、可視の太陽光ではない。

\begin{classroomqa}
\lecturetimestamp{01:03:51}
\audiencequestion{では、我々の望遠鏡は約\(300{,}000\)年より前を見ることができないのですか。}
\lecturerresponse{電磁望遠鏡ではできない。その時代より前、宇宙は電離して不透明であった。他のメッセンジャーは、少なくとも原理的には、より早い時代を探ることができる。}
\end{classroomqa}

\begin{classroomqa}
\lecturetimestamp{01:08:43}
\audiencequestion{光子の温度について話すとき、そのピーク放射を持つ黒体の温度を意味します。しかしそれは波長に対して線形ではないですよね。}
\lecturerresponse{その通り。温度は本質的に典型的な光子のエネルギーであり、したがって光子については波長に反比例してスケールする。\(T^4\)則はエネルギー密度を指し、単一光子のエネルギーではない。それは光子あたりのエネルギーと単位体積あたりの光子数を組み合わせ、数密度は\(T^3\)のようにスケールする。}
\end{classroomqa}

したがって
\begin{equation}
n_\gamma\propto T^3,
\qquad
\rho_\gamma\propto T^4.
\end{equation}

\section{回転曲線から見た暗黒物質}

暗黒物質は銀河の回転曲線を通じて入ってくる。それは大規模なフリードマン方程式から推論されたのではなく、銀河スケールの天文学から推論された。

最も単純なニュートン的問題から始める。質量\(m\)の軽い天体が半径\(r\)で中心質量\(M\)の周りを回っている。重力は
\begin{equation}
\frac{GMm}{r^2},
\end{equation}
であり、向心加速度は\(v^2/r\)である。ニュートンの法則により
\begin{equation}
m\frac{v^2}{r}=\frac{GMm}{r^2}.
\end{equation}
\(m\)を消去し、\(r\)を掛けると
\begin{equation}
v(r)=\sqrt{\frac{GM}{r}}.
\end{equation}
が得られる。したがって、重力質量が中心付近に集中しているなら、軌道速度は\(r^{-1/2}\)のように減少する。

これは質量が光に従うなら期待されることである。銀河の可視物質のほとんどは核の近くにあり、したがって外側の恒星は太陽系外縁の惑星とよく似た振る舞いをするであろう。天文学者は、近くの銀河全体の後退運動を差し引いた後、我々に近づく側と遠ざかる側を比較して、銀河を横切るドップラーシフトの差を測定することにより軌道速度を推論する。測定された軌道速度は大きな半径までほぼ一定にとどまる。

そこで中心質量\(M\)を囲まれた質量\(M(r)\)に置き換える。
\begin{equation}
\frac{v^2}{r}=\frac{G M(r)}{r^2}.
\end{equation}
もし\(v(r)\)がほぼ一定なら、
\begin{equation}
M(r)\propto r.
\end{equation}
言い換えれば、重力質量は可視光がそうでないところでも半径とともに増え続ける。

ここで、議論が単純化のためにそうするように、質量分布がほぼ球対称であると仮定する。すると
\begin{equation}
M(R)=\int_0^R 4\pi r^2\rho(r)\,dr.
\end{equation}
もし\(M(R)\propto R\)なら、\(R\)に関して微分すると
\begin{equation}
4\pi R^2\rho(R)\propto 1,
\end{equation}
となり、したがって
\begin{equation}
\rho(r)\propto \frac{1}{r^2}.
\end{equation}

\begin{classroomqa}
\lecturetimestamp{01:28:18}
\audiencequestion{密度は\(1/r^2\)のように振る舞うべきに見えます。それには理由があるのですか。}
\lecturerresponse{観測された平坦な回転曲線から\(M(R)\propto R\)が推論され、微分すると\(\rho(r)\propto 1/r^2\)が得られる。それ以上については、ここでは単純な第一原理からの説明は主張されない。詳細なシミュレーションが多くの仕事をする。}
\end{classroomqa}

安定した結論は、銀河には、光る物質だけから示唆されるよりもはるかに多くの重力を及ぼす質量が含まれており、その大きさは数倍程度、ここで用いた見積もりではほぼ一桁程度大きいということである。可視の銀河は、はるかに大きな暗黒ハローの中に位置している。

普通の気体がそのような希薄な状態にとどまらないのはなぜか。散逸である。ガス雲は、衝突、摩擦に似た過程、放射によってエネルギーを失えるため収縮する。暗黒物質は明らかにほとんど放射せず、ほぼ無衝突である。重力はそれを内側へ引き寄せ、位置エネルギーを運動エネルギーに変換できるが、効率的な散逸がなければ、普通の物質が作るようなコンパクトな発光構造へと崩壊しない。ハローを広がったままにしているのは運動エネルギーである。

次にいくつかの代替案が慎重に論じられる。ニュートリノは銀河スケールでこのように振る舞うには速すぎると言われる。普通の水素からなるハローは初期宇宙の元素合成をひどく損なうであろう。普通の物質からなるブラックホールも問題があると述べられる。なぜなら、それらは初期宇宙にあまりに多くの普通の物質を意味することになり、レンズ効果の問題を引き起こすからである。その場の議論はそれ以上を主張しない。

ハローは薄い殻やトーラスとして記述されない。それはほぼ球形で一様でない球であり、中心に向かってより高密度で、外側へ徐々に減少する。

\begin{classroomqa}
\lecturetimestamp{01:53:12}
\audiencequestion{弾丸銀河団について話しているのですか。}
\lecturerresponse{はい。X線ガスは普通の物質を示し、レンズ効果は重力を及ぼす全質量を示す。衝突ではそれら二つの分布が分離する。これは、可視物質だけに結びついた単なる重力の修正ではなく、実在の暗黒物質の強い証拠である。}
\end{classroomqa}

これが最後の観測的論点である。暗黒物質が何であるかについて詳細にはほとんど知られていないが、ここでは根本的に新しい原理を必要としない。保守的な可能性は、自然がまだ同定されていない弱く相互作用する粒子を追加で含むというだけである。
